% !TEX root = ../../ctfp-print.tex

\lettrine[lhang=0.17]{现}{在我们已经知道}单子(monad)的作用了——它允许我们组合装饰(embellished)函数——真正有趣的问题是为什么装饰函数在函数式编程中如此重要。我们已经见过一个例子，即\code{Writer}单子，其中装饰使我们能够创建并累积跨越多个函数调用的日志。这个问题原本需要用不纯的函数来解决（例如通过访问和修改某些全局状态），但现在我们用纯函数解决了。

\section{问题}

以下是类似问题的一个简短列表，复制自\urlref{https://core.ac.uk/download/pdf/21173011.pdf}{Eugenio Moggi的开创性论文}，这些问题传统上都通过放弃函数纯洁性来解决：

\begin{itemize}
  \tightlist
  \item
        部分性(Partiality)：可能不会终止的计算
  \item
        非确定性(Nondeterminism)：可能返回多个结果的计算
  \item
        副作用(Side effects)：访问/修改状态的计算

        \begin{itemize}
          \tightlist
          \item
                只读状态(read-only state)，或称环境(environment)
          \item
                只写状态(write-only state)，或称日志(log)
          \item
                读写状态(read/write state)
        \end{itemize}
  \item
        异常(Exceptions)：可能失败的部分函数
  \item
        续体(Continuations)：保存程序状态并在需要时恢复的能力
  \item
        交互式输入(Interactive Input)
  \item
        交互式输出(Interactive Output)
\end{itemize}

真正令人惊叹的是，所有这些问题都可以用同一个巧妙技巧解决：转向装饰函数。当然，每种情况下的装饰形式会完全不同。

你必须认识到，在这个阶段我们并不要求装饰必须是单子性的(monadic)。只有当我们坚持组合操作——即将单个装饰函数分解为更小的装饰函数——时，才需要单子结构。同样，由于每种装饰形式都不同，单子组合的具体实现也会不同，但总体模式是相同的。这是一个非常简单的模式：满足结合律并具有单位元的组合操作。

下一节包含大量Haskell示例。如果你急于回到范畴论内容或已熟悉Haskell的单子实现，可自由略读甚至跳过本节。

\section{解决方案}

首先，让我们分析一下我们使用\code{Writer}单子的方式。我们从一个执行特定任务的纯函数开始——给定参数，它会产生特定的输出。我们用另一个函数替换了这个函数，该函数通过将原始输出与字符串配对来增强原始输出。这是我们解决日志问题的方案。

但我们不能就此止步，因为一般来说，我们不想处理整体式的解决方案。我们需要能够将一个产生日志的函数分解为更小的产生日志的函数。正是这些较小函数的组合引导我们得出了单子的概念。

真正令人惊叹的是，这种增强函数返回类型模式适用于大量通常需要放弃纯度的问题。让我们回顾我们的列表，并依次确定适用于每个问题的增强方式。

\subsection{偏函数性}

我们通过将可能不会终止的每个函数的返回类型转换为"提升"类型来修改其返回类型——包含原始类型所有值以及特殊"底值"$\bot$的类型。例如，作为集合的\code{Bool}类型包含两个元素：\code{True}和\code{False}。提升后的\code{Bool}包含三个元素。返回提升\code{Bool}的函数可能产生\code{True}或\code{False}，或者无限执行下去。

有趣的是，在像Haskell这样的惰性语言中，永不停止的函数实际上可能返回一个值，这个值可以传递给下一个函数。我们将这个特殊值称为底值。只要这个值不是显式需要的（例如用于模式匹配或生成输出），就可以在不阻塞程序执行的情况下传递它。因为每个Haskell函数都可能是非终止的，所以Haskell中的所有类型都被认为是提升的。这就是为什么我们经常谈论Haskell的范畴$\Hask$而不是简单的$\Set$。不过，目前尚不清楚$\Hask$是否是一个真正的范畴（参见这篇\urlref{http://math.andrej.com/2016/08/06/hask-is-not-a-category/}{Andrej Bauer的文章}）。

\subsection{不确定性}

如果一个函数可以返回许多不同结果，不如让它一次性返回所有结果。从语义上看，不确定函数相当于返回结果列表的函数。这在惰性垃圾回收语言中很有意义。例如，如果只需要一个值，可以只取列表的头部，尾部永远不会被求值。如果需要随机值，可以用随机数生成器选择列表的第n个元素。惰性甚至允许你返回无限的结果列表。

在列表单子——Haskell对不确定计算的实现中——\code{join}被实现为\code{concat}。记住\code{join}应该展平容器的容器——\code{concat}将列表的列表连接成单个列表。\code{return}创建一个单元素列表：

\src{snippet01}
列表单子的绑定操作符由通用公式给出：\code{fmap}后接\code{join}，在此情况下给出：

\src{snippet02}
这里，函数\code{k}（本身产生列表）应用于列表\code{as}的每个元素。结果是一个列表的列表，使用\code{concat}将其展平。

从程序员的角度来看，使用列表比调用不确定函数的循环更容易，或者实现返回迭代器的函数（尽管，在\urlref{http://ericniebler.com/2014/04/27/range-comprehensions/}{现代C++}中，返回惰性范围几乎等同于返回Haskell的列表）。

创造性使用不确定性的典型案例出现在游戏编程中。例如，当计算机与人类下国际象棋时，它无法预测对手的下一步行动。但它可以生成所有可能移动的列表并逐一分析。类似地，不确定解析器可以生成给定表达式的所有可能解析结果。

尽管我们可以将返回列表的函数解释为不确定的，但列表单子的应用远不止于此。这是因为将生成列表的计算拼接在一起是命令式编程中迭代结构（循环）的完美函数式替代品。单个循环通常可以用\code{fmap}重写，将循环体应用于列表的每个元素。列表单子的\code{do}表示法可用于替换复杂的嵌套循环。

我最喜欢的案例是生成毕达哥拉斯三元组的程序——能构成直角三角形边长的正整数三元组。

\src{snippet03}
第一行告诉我们\code{z}从无限正整数列表\code{{[}1..{]}}中获取元素。然后\code{x}从有限列表\code{{[}1..z{]}}中获取元素。最后\code{y}从\code{x}到\code{z}的数字列表中获取元素。现在我们有三个满足$1 \leqslant x \leqslant y \leqslant z$的数字。函数\code{guard}接受一个布尔表达式并返回单位列表：

\src{snippet04}
这个函数（属于名为\code{MonadPlus}的更大类）在这里用于过滤非毕达哥拉斯三元组。事实上，如果你查看绑定（或相关操作符\code{>>}）的实现，会注意到当给定空列表时，它会产生空列表。另一方面，当给定非空列表（此处是包含单位的单元素列表\code{{[}(){]}}）时，绑定会调用延续\code{return (x, y, z)}，生成经过验证的毕达哥拉斯三元组的单元素列表。所有这些单元素列表将通过外层绑定连接起来，产生最终（无限）结果。当然，\code{triples}的调用者永远无法消费整个列表，但这无关紧要，因为Haskell是惰性的。

原本需要三层嵌套循环的问题，在列表单子和\code{do}表示法的帮助下得到了戏剧性简化。不仅如此，Haskell还允许你进一步简化代码，使用列表推导式：

\src{snippet05}
这只是列表单子（严格来说是\code{MonadPlus}）的进一步语法糖。

你可能会在其他函数式或命令式语言中看到类似的构造，以生成器和协程的形式存在。

\subsection{只读状态}

具有对外部状态（或环境）只读访问权限的函数，总可以用额外接收该环境作为参数的函数代替。纯函数\code{(a, e) -> b}（其中\code{e}是环境类型）乍看之下不像克莱斯利箭头。但一旦柯里化为\code{a -> (e -> b)}，我们就能识别出增强部分是我们熟悉的读者函子：

\src{snippet06}
你可以将返回\code{Reader}的函数解释为生成微型可执行文件：给定环境产生所需结果的操作。有一个辅助函数\code{runReader}来执行此类操作：

\src{snippet07}
它可能因环境的不同值产生不同的结果。

请注意，返回\code{Reader}的函数和\code{Reader}操作本身都是纯的。

要实现\code{Reader}单子的绑定，首先注意到你需要生成一个接收环境\code{e}并产生\code{b}的函数：

\src{snippet08}
在lambda内部，我们可以执行操作\code{ra}来生成\code{a}：

\src{snippet09}
然后我们可以将\code{a}传递给延续\code{k}以获得新操作\code{rb}：

\src{snippet10}
最后，我们可以用环境\code{e}运行操作\code{rb}：

\src{snippet11}
要实现\code{return}，我们创建一个忽略环境并返回未更改值的操作。

经过一些简化后，我们得到以下定义：

\src{snippet12}

\subsection{只写状态}

这正是我们最初的日志示例。增强部分由\code{Writer}函子给出：

\src{snippet13}
为了完整性，还有一个简单的辅助函数\code{runWriter}用于解包数据构造：

\src{snippet14}
正如我们之前所见，为了让\code{Writer}可组合，\code{w}必须是幺半群。以下是用绑定操作符编写的\code{Writer}单子实例：

\src{snippet15}

\subsection{状态}

具有读/写状态访问权限的函数结合了\code{Reader}和\code{Writer}的增强。你可以将它们视为纯函数，接收状态作为额外参数并产生值/状态对作为结果：\code{(a, s) -> (b, s)}。柯里化后，我们得到克莱斯利箭头形式\code{a -> (s -> (b, s))}，其增强抽象在\code{State}函子中：

\src{snippet16}
同样，我们可以将克莱斯利箭头视为返回操作，可以通过辅助函数执行：

\src{snippet17}
不同的初始状态不仅可能产生不同的结果，还可能产生不同的最终状态。

\code{State}单子的绑定实现与\code{Reader}单子非常相似，只是需要注意在每一步传递正确的状态：

\src{snippet18}
这是完整的实例：

\src{snippet19}
还有两个辅助克莱斯利箭头可用于操作状态。其中一个检索状态供检查：

\src{snippet20}
另一个则用全新的状态替换它：

\src{snippet21}

\subsection{异常}

抛出异常的命令式函数本质上是一个偏函数——它是对某些参数值未定义的函数。用纯全函数实现异常最简单的做法是使用\code{Maybe}函子。偏函数扩展为全函数，当有意义时返回\code{Just a}，无意义时返回\code{Nothing}。如果我们还想返回失败原因的信息，可以改用\code{Either}函子（第一个类型固定，例如\code{String}）。

以下是\code{Maybe}的\code{Monad}实例：

\src{snippet22}
注意当检测到错误时，\code{Maybe}的单子组合会正确短路计算（不会调用延续\code{k}）。这正是我们期望的异常行为。

\subsection{延续}

这是你在求职面试后可能经历的"别打电话给我们，我们会打给你！"的情况。你没有直接得到答案，而是需要提供一个处理程序——一个将在结果可用时被调用的函数。这种编程风格在调用时结果未知的情况下特别有用，例如结果由另一个线程计算或从远程网站交付。在这种情况下，克莱斯利箭头返回一个接受处理程序（代表"计算的剩余部分"）的函数：

\src{snippet23}
当处理程序\code{a -> r}最终被调用时，它产生类型为\code{r}的结果，并在最后返回。延续由结果类型参数化。（实践中，这通常是某种状态指示器。）

还有一个辅助函数用于执行克莱斯利箭头返回的操作。它接收处理程序并将其传递给延续：

\src{snippet24}
延续的组合 notoriously 复杂，因此通过单子特别是\code{do}表示法来处理具有极大优势。

让我们推导绑定的实现。首先看简化签名：

\src{snippet25}
我们的目标是创建一个接收处理程序\code{(b -> r)}并产生结果\code{r}的函数。这就是我们的起点：

\src{snippet26}
在lambda内部，我们想用代表计算剩余部分的处理程序调用函数\code{ka}。我们将这个处理程序实现为lambda：

\src{snippet27}
在这种情况下，计算的剩余部分涉及先用\code{a}调用\code{kab}，然后将\code{hb}传递给生成的操作\code{kb}：

\src{snippet28}
如你所见，延续是内外组合的。最终处理程序\code{hb}从计算的最内层调用。这是完整实例：

\src{snippet29}

\subsection{交互式输入}

这是最棘手的问题也是很多困惑的根源。显然，像\code{getChar}这样的函数如果返回键盘输入的字符，就不可能是纯的。但如果它在容器中返回字符呢？只要无法从容器中提取字符，我们就可以声称该函数是纯的。每次调用\code{getChar}都会返回完全相同的容器。概念上，这个容器包含所有可能字符的叠加态。

如果你熟悉量子力学，应该很容易理解这个类比。就像装着薛定谔猫的盒子——只是无法打开或窥视内部。这个盒子由特殊的内置\code{IO}函子定义。在我们的示例中，\code{getChar}可以声明为克莱斯利箭头：

\src{snippet30}
（实际上，由于从单位类型的函数等价于选择返回类型的值，\code{getChar}的声明简化为\code{getChar :: IO Char}。）

作为函子，\code{IO}允许你使用\code{fmap}操作其内容。而且，作为函子，它可以存储任何类型的值，不仅仅是字符。这种方法的真正实用之处在于，在Haskell中，\code{IO}是一个单子。这意味着你可以组合生成\code{IO}对象的克莱斯利箭头。

你可能认为克莱斯利组合会允许你窥视\code{IO}对象的内容（继续量子类比的话就是"波函数坍缩"）。确实，你可以将\code{getChar}与另一个接收字符并将其转换为整数的克莱斯利箭头组合。关键在于第二个克莱斯利箭头只能返回\code{IO Int}。再次，你会得到所有可能整数的叠加态。依此类推。薛定谔的猫永远不会出袋子。一旦进入\code{IO}单子，就无法逃脱。没有类似\code{runState}或\code{runReader}的\code{IO}单子方法。没有\code{runIO}！

那么除了与其他克莱斯利箭头组合外，还能对克莱斯利箭头的结果——即\code{IO}对象做什么呢？好吧，你可以从\code{main}返回它。在Haskell中，\code{main}的签名为：

\src{snippet31}
你可以自由地将其视为一个克莱斯利箭头：

\src{snippet32}
从这个角度看，Haskell程序只是一个大的\code{IO}单子克莱斯利箭头。你可以使用单子组合从更小的克莱斯利箭头构建它。运行时系统负责对生成的\code{IO}对象（也称为\code{IO}操作）做某些事情。

请注意，箭头本身是纯函数——纯函数贯穿始终。脏活累活被推给了系统。当它最终执行从\code{main}返回的\code{IO}操作时，它会做各种讨厌的事情，如读取用户输入、修改文件、打印烦人消息、格式化磁盘等等。Haskell程序永远不会弄脏自己的手（嗯，除非调用\code{unsafePerformIO}，但那是另一个故事）。

当然，因为Haskell是惰性的，\code{main}几乎立即返回，而脏活立即开始。在执行\code{IO}操作期间，按需请求和求值纯计算结果。因此，实际上，程序的执行是纯（Haskell）代码和脏（系统）代码的交错。

有一种更离奇但作为数学模型完全合理的\code{IO}单子解释。它将整个宇宙视为程序中的一个对象。概念上，命令式模型将宇宙视为外部全局对象，因此执行I/O的过程通过与该对象交互产生副作用。它们既可以读取也可以修改宇宙的状态。

我们已经知道如何在函数式编程中处理状态——我们使用状态单子。然而，与简单状态不同，宇宙状态无法轻易用标准数据结构描述。但只要我们不直接与其交互，就不需要这么做。只要假设存在一个类型\code{RealWorld}，并且通过某种宇宙工程奇迹，运行时能够提供该类型的对象即可。一个\code{IO}操作只是一个函数：

\src{snippet33}
或者，用状态单子表示：

\src{snippet34}
然而，\code{IO}单子的\code{>=>}和\code{return}必须内置于语言中。

\subsection{交互式输出}

同一个\code{IO}单子也用于封装交互式输出。\code{RealWorld}应该包含所有输出设备。你可能想知道为什么我们不能直接从Haskell调用输出函数并假装它们什么都不做。例如，为什么我们要：

\src{snippet35}
而不是更简单的：

\src{snippet36}
两个原因：Haskell是惰性的，所以它永远不会调用输出（这里是单位对象）未被使用的函数。即使它不是惰性的，它仍然可以自由改变这些调用的顺序，从而破坏输出。在Haskell中强制两个函数顺序执行的唯一方法是通过数据依赖。一个函数的输入必须依赖于另一个函数的输出。在\code{IO}操作之间传递\code{RealWorld}强制序列化。

概念上，在这个程序中：

\src{snippet37}
打印"World!"的操作接收作为输入的"Hello "已显示在屏幕上的宇宙。它输出一个新的宇宙，屏幕上显示"Hello World!"。

\section{结论}

当然，我仅仅触及了单子编程(monadic programming)的表面。
单子不仅通过纯函数(pure functions)实现了通常由命令式编程(imperative programming)中的副作用(side effects)完成的功能，
而且实现了更高程度的控制和类型安全性(type safety)。
不过，单子也并非没有缺点。
关于单子的主要批评在于它们彼此之间不易组合。
诚然，你可以使用单子变换器库(mond transformer library)来组合大多数基本单子。
创建一个结合状态(state)与异常(exceptions)的单子栈(mond stack)相对容易，
但并没有通用公式能将任意单子组合在一起。
